\subsection{Why the local exceptional ranks give Smith tails}\label{sec:lattice-walk}

The step between a joint rank estimate and Proposition~\ref{prop:local-index} is a finite lattice argument. The following is the source's \texttt{WalkSharp}, specialized to its actual parameters; the proof explains both the exponent $m^{\beta+1}$ and the depth $\tau+\beta S$. \textsf{[LW1--LW3]}

Work over a DVR $\mathcal O$ containing $w$ with $w^2=p$. Order the local indices $1,\ldots,V$. There is a compatible lattice $L_\sigma$ for each choice of a base, ball, or dual slot at every index. Write $M=L_{\rm base}$ and $L_B$ for dual slots on $B$, base slots elsewhere. With the other slots fixed, assume
\[
 L_{\rm dual}\subseteq L_{\rm base}\subseteq L_{\rm ball},\qquad
 w^{t_v}L_{\rm ball}\subseteq L_{\rm base},\qquad
 w^{t_v}L_{\rm base}\subseteq L_{\rm dual},\qquad
 y_vL_{\rm ball}\subseteq L_{\rm dual}.
\]
The $y_v$ commute. In particular they preserve $M$. Suppose, for every nonempty $B$, that a finite set $G_B\subseteq L_B$ satisfies
\begin{equation}\label{eq:walk-joint}
 p^{D_0+1}L_B\subseteq
 \sum_{v\in B}y_vL_{B;v=\rm ball}+\langle G_B\rangle,
 \qquad |G_B|\le d\prod_{v\in B}\kappa_v,
\end{equation}
where $d=\operatorname{rank}M$ and $\kappa_v\ge0$. The preceding joint-rank argument supplies this estimate from the single-place inputs; it is not a statement about arbitrary intersecting subspaces.

For integers $m\ge1$ and $t_v,D_0,S_v,S,\beta\ge0$, put $h_v=t_v+D_0+1$, $\tau=m\sum_vh_v$, and
\[
 N=\sum_{A\subseteq\{1,\ldots,V\}}
 p^{\sum_{v\notin A}S_v}\Bigl(\prod_{v\in A}y_v^m\Bigr)M,
 \qquad 0\le S_v\le S.
\]
Then $N\subseteq M$, $p^{\sum_vS_v}M\subseteq N$, and for every $\beta\ge0$,
\begin{equation}\label{eq:walk-tail}
 p^{\tau+\beta S}M\subseteq N+\langle C_\beta\rangle,
 \qquad |C_\beta|\le d\,m^{\beta+1}e_{\beta+1}(\kappa).
\end{equation}
For an integral count one may take the floor of the right-hand side. This is precisely the exceptional-direction hypothesis used below.

\begin{proof}
Set $q=\beta+1$. If $q>V$, then $\sum S_v\le\beta S$ and the assertion follows with $C_\beta=\varnothing$. Otherwise keep $q$ active indices. A state consists of a set $B$ of size $q$ and counters $0\le i_v<m$ on $B$. The indices less than $\max B$ and outside $B$ are completed; larger indices are waiting. Define
\[
 Y_{B,i}=\prod_{u\,\rm completed}y_u^m\prod_{v\in B}y_v^{i_v},
 \quad
 A_{B,i}=\sum_{u\,\rm completed}2mh_u+
                 \sum_{v\in B}(t_v+2i_vh_v).
\]
Use as $C_\beta$ the union of all $Y_{B,i}G_B$. Every vector belongs to $M$, and each $B$ has exactly $m^q$ choices of counters. Hence its cardinality is at most
$m^q\sum_{|B|=q}|G_B|\le d m^q e_q(\kappa)$.

Let $E=2(\tau+\beta S)$. We prove the state invariant
\[
 w^{E-A_{B,i}}Y_{B,i}L_B\subseteq N+\langle C_\beta\rangle
\]
backwards in the sum of the operator exponents, which strictly increases at every transition and is bounded by $mV$. Apply \eqref{eq:walk-joint}, spending $2(D_0+1)$ powers of $w$. Its exceptional term is already in $\langle C_\beta\rangle$. Each other term has the form $Y_{B,i}y_vz$ with $z\in L_{B;v=\rm ball}$.

If $i_v+1<m$, multiply $z$ by $w^{2t_v}$ to return from ball to dual, and increase $i_v$ by one. The budget increases by $2h_v$, exactly the sum of these two costs. If $i_v+1=m$ and an index is waiting, return the completed $v$ from ball to base at cost $t_v$, enter the next waiting index $u$ in its dual slot at cost $t_u$, and replace $v$ by $u$ in $B$. The budget increases by
$2(D_0+1)+t_v+t_u$, again exactly the amount spent.

Finally suppose $i_v+1=m$ and no index is waiting. All but the remaining $\beta$ active indices now carry $y_u^m$. After paying $t_v$ to return to base, pay
$2\sum_{u\in B\setminus\{v\}}S_u\le2\beta S$ to use the corresponding generator of $N$. Each completed index has spent $2mh_u$; an active index has spent $t_u+2i_uh_u\le2mh_u$, with enough unused budget for the next decomposition or final exit. Thus all the displayed powers of $w$ are nonnegative, and the backwards induction is valid.

Start with $B=\{1,\ldots,q\}$ and all counters zero. Entering its dual slots costs $\sum_{v\in B}t_v=A_{B,0}$, so the invariant gives $w^EM\subseteq N+\langle C_\beta\rangle$. Since $w^2=p$, this is \eqref{eq:walk-tail}.
\end{proof}

The three slots themselves have a simple lattice interpretation. Locally the ball is
$w^{-t}H\cap\bigcap_eR_e$, and the dual slot is the integral polar of
$w^{-t}H^\vee\cap\bigcap_eR'_e$. Once $H\subseteq R_e$ and $H^\vee\subseteq R'_e$ are known, reflexivity gives the four displayed lattice inclusions. Tensoring and taking integral equivariant functions preserves them. The edge inclusions and the compatibility of the adjacency operators with the polar pairing remain explicit source inputs here; their proofs use the coefficient representations. \textsf{[LW4]}
